\chapter{Optimisation et calcul numérique}
\label{ch:optim}

\begin{lecture}
La fonction de coût associe un nombre à chaque choix des coefficients. Optimiser
consiste à déplacer progressivement ces coefficients vers une zone où ce nombre
est plus faible. Le gradient donne la direction de plus forte hausse ; la
descente suit donc la direction opposée, avec une longueur de pas réglée par
$\alpha$.
\end{lecture}

La descente de gradient simultanée produit la suite
\[
  \thv^{[t+1]}=\thv^{[t]}-\alpha\nabla J(\thv^{[t]}),
  \qquad \alpha>0.
\]
Cette itération est la descente de gradient par lot appliquée à la
log-vraisemblance négative \cite[chap.~8]{boyd2004}.

Le gradient indique la direction locale de hausse du coût ; son opposé indique
donc une direction locale de baisse. Cette propriété résulte du développement
au premier ordre présenté à l'annexe~\ref{app:descent-direction}. Elle reste
locale : un pas trop grand peut dépasser la zone où cette approximation décrit
correctement la fonction.
Le mot \emph{simultanée} est important : toutes les composantes du nouveau
vecteur sont calculées à partir du même ancien vecteur. Un taux trop grand peut
faire augmenter le coût ou provoquer une divergence ; un taux très petit rend
la convergence inutilement lente.

\begin{figure}[htbp]
  \centering
  \begin{tikzpicture}
    \begin{axis}[courseplot,axis lines=left,
      xlabel={itération $t$},ylabel={$J(\thv^{[t]})$},xmin=0,xmax=30,
      ymin=.18,ymax=1.05,grid=major,grid style={gridgray!50},samples=120]
      \addplot[accent,very thick,domain=0:30]{.22+.75*exp(-x/5)};
      \addplot[warning,dashed,thick,domain=0:30]{.32+.55*exp(-x/12)+.08*sin(deg(x))};
    \end{axis}
  \end{tikzpicture}
  \caption{Diagnostic typique du coût : décroissance stable (trait plein) et
  oscillations dues à un pas trop agressif (pointillé). Cette figure est
  schématique ; le tracé du coût réel doit être conservé avec l'expérience.}
\end{figure}

Un arrêt robuste combine un nombre maximal d'itérations avec, par exemple,
\[
  \frac{|J^{(t+1)}-J^{(t)}|}{\max(1,|J^{(t)}|)}<\eps
  \quad\text{ou}\quad \|\nabla J(\thv^{[t]})\|_2<\eps.
\]
Un simple nombre fixe d'itérations ne démontre pas la convergence.

\begin{figure}[htbp]
\centering
\begin{tikzpicture}
\begin{axis}[width=.92\textwidth,height=11.2cm,
  xlabel={$\theta_0$},ylabel={$\theta_1$},xmin=-3,xmax=3,ymin=-3,ymax=3,
  axis equal image,grid=major,grid style={gridgray!30}]
  \foreach \a/\b in {.35/.65,.65/1.2,1/1.85,1.35/2.5}{
    \addplot[accent!55,thick,domain=0:360,samples=140]
      ({-.30+\a*cos(x)},{.35+\b*sin(x)});
  }
  \addplot[warning,very thick,mark=*,mark size=2.1pt,-{Latex}]
    coordinates {(-2.55,2.55) (-1.65,1.92) (-1.05,1.32) (-.68,.86)
      (-.46,.56) (-.34,.40) (-.30,.35)};
  \addplot[only marks,mark=star,mark size=5pt,ink]
    coordinates {(-.30,.35)} node[above right]{minimum};
\end{axis}
\end{tikzpicture}
\caption{Courbes de niveau schématiques de $J(\theta_0,\theta_1)$ et trajectoire
de descente. L'allongement des ellipses traduit un mauvais conditionnement : un
même taux d'apprentissage progresse vite dans une direction et lentement dans
l'autre. Les points successifs aboutissent ici au centre des courbes, le minimum.
La standardisation tend à réduire cet effet d'allongement.}
\label{fig:contours}
\end{figure}

Il faut distinguer trois diagnostics. Une baisse du coût indique que l'algorithme
avance ; une petite norme du gradient indique qu'il approche un point
stationnaire ; une bonne performance de validation indique que ce point produit
un modèle utile hors échantillon. Aucun de ces faits n'implique automatiquement
les deux autres.

Un \emph{point stationnaire} est un point où le gradient est nul. Pour
$J(\theta)=\theta^2$, le point $0$ est stationnaire et constitue le minimum. Pour
$J(\theta)=\theta^3$, le point $0$ est aussi stationnaire, mais la fonction
continue de décroître à gauche et de croître à droite : ce n'est pas un minimum.
Dans le problème logistique convexe, un point stationnaire fini est un minimum
global ; c'est la convexité établie au chapitre précédent qui autorise cette
conclusion, pas la seule égalité $\nabla J=0$.

\paragraph{Influence du taux d'apprentissage dans le cas scalaire.}
La fonction $J(\theta)=(\theta-2)^2$ mesure ici la distance quadratique au
minimum situé en $\theta=2$. À chaque tour, le programme évalue la pente
$J'(\theta)=2(\theta-2)$, la multiplie par $\alpha$, puis soustrait cette
correction à la valeur courante. Depuis $\theta^{[0]}=-2$, le pas
$\alpha=0{,}25$ rapproche successivement le paramètre du minimum, tandis que
$\alpha=1$ le fait passer alternativement d'un côté à l'autre sans réduire
l'écart.

\begin{figure}[htbp]
\centering
\begin{minipage}[t]{.48\textwidth}
\centering
\begin{tikzpicture}
\begin{axis}[width=.85\linewidth,height=6.2cm,axis lines=left,
  title={$\alpha=0{,}25$ : rapprochement du minimum},
  xlabel={itération $t$},ylabel={paramètre $\theta^{[t]}$},
  xmin=0,xmax=6,ymin=-2.6,ymax=6.6,xtick={0,1,2,3,4,5,6},
  grid=major,grid style={gridgray!45}]
  \addplot[accent,very thick,mark=*] coordinates
    {(0,-2) (1,0) (2,1) (3,1.5) (4,1.75) (5,1.875) (6,1.9375)};
  \addplot[ink,dashed,thick,domain=0:6]{2};
  \node[anchor=south east] at (axis cs:6,2){minimum $\theta=2$};
\end{axis}
\end{tikzpicture}
\end{minipage}\hfill
\begin{minipage}[t]{.48\textwidth}
\centering
\begin{tikzpicture}
\begin{axis}[width=.85\linewidth,height=6.2cm,axis lines=left,
  title={$\alpha=1$ : oscillation persistante},
  xlabel={itération $t$},ylabel={paramètre $\theta^{[t]}$},
  xmin=0,xmax=6,ymin=-2.6,ymax=6.6,xtick={0,1,2,3,4,5,6},
  grid=major,grid style={gridgray!45}]
  \addplot[warning,very thick,mark=*] coordinates
    {(0,-2) (1,6) (2,-2) (3,6) (4,-2) (5,6) (6,-2)};
  \addplot[ink,dashed,thick,domain=0:6]{2};
  \node[anchor=south east] at (axis cs:6,2){minimum $\theta=2$};
\end{axis}
\end{tikzpicture}
\end{minipage}
\caption{Valeur du paramètre après chaque mise à jour. Les segments relient deux
itérations successives ; ils représentent l'ordre des calculs, pas une grandeur
supplémentaire.}
\label{fig:learning-rate-scalar}
\end{figure}

Dans le modèle complet, la même séquence s'applique à un vecteur : calculer les
scores, les probabilités, le gradient, puis mettre à jour simultanément tous les
coefficients. Le suivi de $J$ permet de détecter une décroissance, une oscillation
ou une divergence.

\begin{figure}[htbp]
\centering
\resizebox{\textwidth}{!}{%
\begin{tikzpicture}[
  b/.style={draw=accent,fill=accentlight,rounded corners=2pt,
    minimum width=27mm,minimum height=11mm,align=center},
  a/.style={-{Latex},thick,draw=ink},node distance=8mm]
  \node[b] (input) {$\X,\y,\thv^{[t]}$\\données et paramètres};
  \node[b,right=of input] (score) {$\mathbf z=\X\thv^{[t]}$\\scores};
  \node[b,right=of score] (prob) {$\mathbf p=\sigma(\mathbf z)$\\probabilités};
  \node[b,right=of prob] (res) {$\mathbf r=\mathbf p-\y$\\écarts};
  \node[b,right=of res] (grad) {$\mathbf g=\X^\top\mathbf r/m$\\gradient};
  \node[b,below=13mm of grad] (update)
    {$\thv^{[t+1]}=\thv^{[t]}-\alpha\mathbf g$\\mise à jour simultanée};
  \node[b,left=of update] (control) {$J,\|\mathbf g\|_2$\\contrôles d'arrêt};
  \draw[a] (input)--(score); \draw[a] (score)--(prob);
  \draw[a] (prob)--(res); \draw[a] (res)--(grad);
  \draw[a] (grad)--(update); \draw[a] (update)--(control);
  \draw[a] (control.west)-| node[pos=.25,below]{itération suivante} (input.south);
\end{tikzpicture}}
\caption{Cycle de calcul d'une itération de descente de gradient par lot. Chaque
flèche correspond à une transformation nécessaire ; aucun réajustement des
données préparées n'intervient dans cette boucle.}
\label{fig:training-cycle}
\end{figure}

\begin{procedure}{Exécuter une itération d'ajustement}
\textbf{Entrées :} $\X$, $\y$, les paramètres courants $\thv^{[t]}$ et le taux
$\alpha$. \textbf{Sorties :} $\thv^{[t+1]}$, le coût $J^{[t]}$ et la norme du
gradient $\|\mathbf g^{[t]}\|_2$.
\begin{enumerate}
  \item calculer $\mathbf z=\X\thv^{[t]}$, $\mathbf p=\sigma(\mathbf z)$,
        $J^{[t]}$ et $\mathbf g^{[t]}=\X^\top(\mathbf p-\y)/m$ ;
  \item calculer toutes les composantes de
        $\thv^{[t+1]}=\thv^{[t]}-\alpha\mathbf g^{[t]}$ depuis le même ancien vecteur ;
  \item ajouter $J^{[t]}$ et $\|\mathbf g^{[t]}\|_2$ aux historiques ;
  \item arrêter si la tolérance ou le nombre maximal d'itérations est atteint ;
  \item sinon remplacer $t$ par $t+1$ ; si $J$ augmente durablement ou devient
        non fini, interrompre et revoir $\alpha$ ou la stabilité des calculs.
\end{enumerate}
\end{procedure}

\section{Formules numériquement stables}

\begin{lecture}
Deux formules mathématiquement égales peuvent se comporter différemment sur un
ordinateur, dont les nombres ont une taille et une précision finies. Les formes
stables évitent de calculer des exponentielles immenses, puis de soustraire ou
de logarithmer des valeurs arrondies à $0$ ou $1$.
\end{lecture}

Évaluer $e^{-z}$ sans précaution peut déborder pour un score très négatif. Une
sigmoïde stable utilise deux branches :
\[
  \sigma(z)=
  \begin{cases}
    1/(1+e^{-z}),&z\geq0,\\
    e^z/(1+e^z),&z<0.
  \end{cases}
\]
Il est préférable de calculer directement la perte à partir du score,
\[
  \ell(y,z)=\log(1+e^z)-yz=\operatorname{softplus}(z)-yz,
\]
L'identité résulte directement du remplacement
$p=\sigma(z)$ dans la log-perte et constitue la forme usuelle de la perte
logistique \cite[sec.~7.1]{murphy2022}.
avec une primitive stable de type $\log(1+e^z)$, plutôt que de calculer
$\log(\sigma(z))$ après arrondi de la probabilité à $0$ ou $1$.

\begin{resultat}[Contrôle du gradient]
Avant tout entraînement, comparer le gradient analytique à une différence finie
centrée sur un cas test de faible dimension :
\[
  \frac{\partial J}{\partial\theta_j}\approx
  \frac{J(\thv+h e_j)-J(\thv-h e_j)}{2h}.
\]
Ce test détecte efficacement un signe inversé, un facteur $1/m$ oublié ou un
problème d'orientation des matrices. Il valide la dérivation, pas l'ensemble du
modèle.
\end{resultat}

\begin{figure}[htbp]
\centering
\begin{tikzpicture}
\begin{axis}[courseplot,axis lines=middle,xlabel={$\theta_j$},ylabel={$J$},
  xmin=-1.1,xmax=2.1,ymin=0,ymax=3.8,xtick={-.3,.2,.7},
  xticklabels={$\theta_j-h$,$\theta_j$,$\theta_j+h$},ytick=\empty,
  legend style={at={(.5,1.03)},anchor=south,legend columns=2,draw=none},samples=180]
  \addplot[accent,very thick,domain=-1.1:2.1]{(x-.7)^2+.45};
  \addlegendentry{$J(\theta_j)$}
  \addplot[warning,very thick,domain=-.3:.7]{-x+1.15};
  \addlegendentry{corde entre les deux évaluations}
  \addplot[only marks,mark=*,warning,mark size=2.8pt]
    coordinates {(-.3,1.45) (.7,.45)};
  \addplot[only marks,mark=*,ink,mark size=2.4pt] coordinates {(.2,.70)};
  \draw[dashed,gridgray] (axis cs:-.3,0)--(axis cs:-.3,1.45);
  \draw[dashed,gridgray] (axis cs:.7,0)--(axis cs:.7,.45);
  \node[anchor=west,align=left] at (axis cs:.82,1.15)
    {$h=0{,}5$\\pente $=(0{,}45-1{,}45)/(2h)=-1$};
\end{axis}
\end{tikzpicture}
\caption{Exemple numérique cohérent avec
$J(u)=(u-0{,}7)^2+0{,}45$, $\theta_j=0{,}2$ et $h=0{,}5$. Les deux points de la
corde appartiennent à la courbe. Sa pente vaut $-1$, exactement comme
$J'(0{,}2)=2(0{,}2-0{,}7)=-1$ dans cet exemple quadratique.}
\label{fig:finite-difference}
\end{figure}

\begin{procedure}{Évaluer les formules sans instabilité numérique}
\textbf{Entrée :} un score réel $z$, éventuellement de grande amplitude.
\textbf{Sorties :} une probabilité et une perte finies, calculées sans dépassement
de capacité.
\begin{enumerate}
  \item si $z\geq0$, utiliser $1/(1+e^{-z})$ ; sinon utiliser
        $e^z/(1+e^z)$, afin que l'exponentielle porte sur un nombre négatif ;
  \item calculer la perte depuis $z$ avec une primitive stable de softplus ;
  \item tester des scores extrêmes, par exemple $z=-1000,0,1000$, et vérifier
        l'absence de valeur infinie ou indéfinie ;
  \item sur un cas de faible dimension, comparer gradient analytique et différence
        finie avec une tolérance fixée avant l'entraînement complet.
\end{enumerate}
\end{procedure}
